% Generated by lean-to-paper from paradoxical/CanonicalShape.lean (2026-07-11).
% (C) 2026 Ralf Stephan, in collaboration with Claude Code. CC0 1.0 Universal.
\documentclass[11pt]{amsart}
\usepackage{amsmath,amssymb,amsthm}
\usepackage[hidelinks]{hyperref}
\usepackage{booktabs}

\theoremstyle{plain}
\newtheorem{theorem}{Theorem}[section]
\newtheorem{proposition}[theorem]{Proposition}
\newtheorem{lemma}[theorem]{Lemma}
\newtheorem{corollary}[theorem]{Corollary}
\newtheorem{conjecture}[theorem]{Conjecture}
\theoremstyle{definition}
\newtheorem{definition}[theorem]{Definition}
\theoremstyle{remark}
\newtheorem{remark}[theorem]{Remark}

\newcommand{\N}{\mathbb{N}}
\newcommand{\Par}{\mathcal{P}}        % the paradoxical set
\newcommand{\ParR}[1]{\Par_{#1}}      % the R-circuit slice

% ---------------------------------------------------------------------------
% Terminology table (lean-to-paper; the source of every technical term used).
%
% Lean identifier / concept   | literature term            | source            | notation
% ----------------------------+----------------------------+-------------------+----------
% CC.T                        | Collatz function           | [Roz25, eq. (1)]  | T
% CC.T_iter j n               | iterate of T               | [Roz25, Sec. 1]   | T^j(n)
% CC.X (parity bit)           | parity (odd term count)    | [Roz25, Def. 2.1] | (absorbed)
% pbits j n                   | parity vector              | [Roz25, Def. 2.1] | V_j(n)
% CC.num_odd_steps j n        | number of odd terms        | [Roz25, Def. 2.1] | q_j(n), q
% CC.C j n = 3^q/2^j          | coefficient                | [Roz25, eq. (2)]  | C_j(n), C
% (iterate sequence)          | the sequence Omega_j(n)    | [Roz25, Def. 4.1] | \Omega_j(n)
% RT.IsParadoxical j n        | paradoxical (sequence)     | [Roz25, Def. 1.1] | --
% run-superscript words       | binary word 1^a 0^e ...    | [Roz25, Def. 2.1] | 1^{a}0^{e}
% (block 1^a 0^e)             | circuit                    | [Roz25, App. A]   | (a, e)
%                             |   (extended usage; see footnote in Sec. 1)
% RCircuitSlice R j n         | (descriptive: paradoxical  | --                | \mathcal{P}_R
%                             |  pairs whose parity vector |
%                             |  is a word of R circuits)  |
% shapeOf w                   | (descriptive: canonical    | run-length coding | \sigma(w)
%                             |  circuit decomposition)    | is standard       |
% wordOfShape S               | (descriptive: the word     | [Roz25, Def. 2.1] | 1^{a_1}0^{e_1}...
%                             |  spelled by the blocks)    |  notation         |
% circuitCount j n            | (descriptive: number of    | --                | R_j(n)
%                             |  canonical circuits)       |
% M_T(n) <= n^beta            | maximum excursion bound    | [Roz25, eq. (16)] | M_T
% d_T(n) <= alpha log n       | delay bound                | [Roz25, Conj. 6.2]| d_T
% ---------------------------------------------------------------------------

\begin{document}

\title{Canonical circuit decompositions of parity vectors and the finiteness of paradoxical Collatz sequences}
\author{Ralf Stephan}
\thanks{In collaboration with Claude Code.  All results stated below are
machine-verified in Lean~4; see Appendix~\ref{app:lean}}
\date{July 11, 2026}

\begin{abstract}
For the Collatz function $T$, Rozier and Terracol call a finite sequence
$\Omega_j(n)=(n,T(n),\dots,T^j(n))$ with $n>2$ \emph{paradoxical} when its
coefficient $C_j(n)=3^{q}/2^{j}$ is less than~$1$ and yet $T^j(n)\ge n$.  We
show that the parity vector of any such sequence decomposes canonically into
\emph{circuits} --- blocks of the form $1^{a}0^{e}$ --- so that the paradoxical
set is exhausted by the family $(\ParR{R})_{R\ge 1}$ of its subsets consisting
of the sequences whose parity vector spells exactly $R$ circuits.  We then
prove that the existence of a bound on the length $j$, valid on $\ParR{R}$
\emph{uniformly in $R$}, is equivalent to the finiteness of the paradoxical
set; by Corollary~3.3 of Rozier--Terracol such a bound would therefore imply
the Collatz conjecture.  This makes precise the sense in which the
uniform-over-$R$ strengthening of the per-$R$ finiteness theorem is not merely
sufficient for, but a reformulation of, the finiteness of the paradoxical set.
All results are elementary and machine-verified in Lean~4.
\end{abstract}

\maketitle

\section{Notation and background}\label{sec:background}

We follow the notation of \cite{Roz25} throughout.  Let $T$ be the Collatz
function
\[
T(n)=
\begin{cases}
\dfrac{3n+1}{2}, & n \text{ odd},\\[1ex]
\dfrac{n}{2}, & n \text{ even},
\end{cases}
\]
on the positive integers, and for $j\ge 0$ let
$\Omega_j(n)=\bigl(T^k(n)\bigr)_{k=0}^{j}$ denote the sequence of iterates of
length $j+1$ \cite[Def.~4.1]{Roz25}.  Following
\cite[Def.~2.1]{Roz25}, $q_j(n)$ (or simply $q$) is the number of odd terms
among $n,T(n),\dots,T^{j-1}(n)$; the \emph{coefficient} of $\Omega_j(n)$ is
\[
C_j(n)=\frac{3^{q}}{2^{j}},
\]
and $V_j(n)\in\{0,1\}^j$ is the \emph{parity vector} of length $j$ induced by
$n$, namely $V_j(n)=\bigl(n,T(n),\dots,T^{j-1}(n)\bigr)\bmod 2$.  As in
\cite[Def.~2.1]{Roz25} we write parity vectors as binary words and abbreviate
repetitions by superscripts, so that $1^{2}0^{3}1$ stands for
$\langle 110001\rangle$.

\begin{definition}[{\cite[Def.~1.1]{Roz25}}]\label{def:paradoxical}
For positive integers $j$ and $n$ with $n>2$, the sequence $\Omega_j(n)$ is
\emph{paradoxical} if
\begin{enumerate}
\item[(i)] $C_j(n)<1$, and
\item[(ii)] $T^j(n)\ge n$.
\end{enumerate}
We write
\[
\Par \;=\; \bigl\{(j,n) : \Omega_j(n) \text{ paradoxical},\ n>2\bigr\}
\]
for the \emph{paradoxical set}.
\end{definition}

\begin{remark}\label{rem:predicate}
In the formal development the two conditions (i)--(ii) constitute the
paradoxicality predicate by themselves, and the side condition $n>2$ is
carried separately by every statement below; the conjunction agrees exactly
with Definition~\ref{def:paradoxical}.  (The exclusion of sequences that
merely repeat the trivial cycle, stipulated for convenience in
\cite[Def.~1.1]{Roz25}, is automatic once $n>2$.)  The formal predicate is
defined for all $j\ge 0$, but for $j=0$ condition (i) fails since
$C_0(n)=1$, so nothing is added beyond the positive lengths considered
in \cite{Roz25}.
\end{remark}

A paradoxical sequence with a unique local maximum is, in the cyclic case,
classically called a \emph{$1$-cycle} or \emph{circuit}; its parity vector is
a single block $1^{q}0^{j-q}$ \cite[App.~A]{Roz25}.  We extend the word to
the individual blocks of a general parity vector:\footnote{In
\cite[App.~A]{Roz25} \emph{circuit} refers to the whole (near-)cyclic
sequence whose parity vector is one block $1^{q}0^{j-q}$.  We are not aware
of standard terminology for the individual blocks and adopt \emph{circuit}
for them, as each block records one ascent--descent phase of the sequence.}
a \emph{circuit} is a pair $(a,e)$ of nonnegative integers, spelling the word
$1^{a}0^{e}$, and a finite list $S=\bigl((a_1,e_1),\dots,(a_R,e_R)\bigr)$ of
circuits spells the concatenated word
\[
w(S) \;=\; 1^{a_1}0^{e_1}\,1^{a_2}0^{e_2}\cdots 1^{a_R}0^{e_R},
\]
of length $\sum_{i=1}^{R}(a_i+e_i)$.

\begin{definition}\label{def:slice}
For $R\ge 0$, the \emph{$R$-circuit slice} $\ParR{R}\subseteq\Par$ is the set
of pairs $(j,n)\in\Par$ such that
\[
V_j(n) \;=\; 1^{a_1}0^{e_1}\cdots 1^{a_R}0^{e_R}
\]
for some nonnegative integers $a_1,e_1,\dots,a_R,e_R$ with
$\sum_{i}(a_i+e_i)=j$.
\end{definition}

Blocks in Definition~\ref{def:slice} are allowed to be empty, so the slices
are not disjoint (padding with the circuit $(0,0)$ embeds $\ParR{R}$ into
$\ParR{R+1}$); the results below concern only their union.  We shall also use
the following finiteness criterion from the formal development accompanying
\cite{Roz25}, which is proved there by bounding the first term of a
paradoxical sequence of length $j$ by $2^{j}3^{j}$.

\begin{proposition}\label{prop:lengthbound}
If there is a $J\in\N$ such that $j\le J$ for every $(j,n)\in\Par$, then
$\Par$ is finite.
\end{proposition}

\section{The canonical circuit decomposition}\label{sec:canonical}

Every binary word is spelled by a canonical list of circuits, obtained by
grouping each maximal block of $1$'s with the maximal block of $0$'s that
follows it.

\begin{definition}\label{def:shape}
The \emph{canonical circuit decomposition} $\sigma(w)$ of a finite binary
word $w$ is defined by recursion: $\sigma(\varepsilon)=\varepsilon$ for the
empty word, and for $w\neq\varepsilon$,
\[
\sigma(w) \;=\; \bigl((a,e)\bigr) \frown \sigma(w''),
\]
where $1^{a}$ is the maximal initial block of $1$'s of $w$, $0^{e}$ is the
maximal block of $0$'s following it, and $w''$ is the remaining suffix.  A
word beginning with $0$ thus yields a first circuit with empty ascent,
$a_1=0$.
\end{definition}

(The formal definition operates on words over $\N$, counting every letter
different from $1$ towards the block $0^{e}$; on binary words this is the
description just given.)  Since each recursion step consumes at least one
letter, $\sigma(w)$ is nonempty whenever $w$ is, and the recursion
terminates.

\begin{proposition}\label{prop:roundtrip}
For every finite binary word $w$, the canonical circuits of $w$ spell $w$:
if $\sigma(w)=\bigl((a_1,e_1),\dots,(a_R,e_R)\bigr)$, then
\[
w \;=\; 1^{a_1}0^{e_1}\cdots 1^{a_R}0^{e_R}.
\]
In particular $\sum_{i=1}^{R}(a_i+e_i)$ equals the length of $w$.
\end{proposition}

\begin{proof}[Proof sketch]
Induction along the recursion of Definition~\ref{def:shape}.  Because $w$ is
binary, the maximal initial block of $1$'s and the following maximal block of
non-$1$'s are the constant words $1^{a}$ and $0^{e}$, and $w$ is their
concatenation with the suffix $w''$.  The inductive hypothesis applies to
$w''$ (a suffix of a binary word being binary), and reassembling the three
pieces recovers $w$.  The length statement follows by comparing lengths on
both sides.
\end{proof}

\section{Coverage: every paradoxical sequence lies in its canonical slice}\label{sec:coverage}

The parity vector $V_j(n)$ is a binary word of length $j$, so
Proposition~\ref{prop:roundtrip} applies to it.

\begin{definition}\label{def:circuitcount}
The \emph{canonical circuit count} $R_j(n)$ is the number of circuits of
$\sigma\bigl(V_j(n)\bigr)$.
\end{definition}

For $j\ge 1$ the word $V_j(n)$ is nonempty and hence $R_j(n)\ge 1$.

\begin{lemma}[Coverage]\label{lem:coverage}
If $\Omega_j(n)$ is paradoxical with $n>2$, then $(j,n)\in\ParR{R_j(n)}$.
\end{lemma}

\begin{proof}[Proof sketch]
The canonical decomposition $\sigma\bigl(V_j(n)\bigr)$ consists of $R_j(n)$
circuits by definition, spells $V_j(n)$ by
Proposition~\ref{prop:roundtrip}, and has total length $j$; it therefore
witnesses the membership required by Definition~\ref{def:slice}.
\end{proof}

\begin{corollary}[Exhaustion]\label{cor:exhaustion}
The circuit slices cover the paradoxical set exactly:
\[
\bigcup_{R\ge 0}\ParR{R} \;=\; \Par .
\]
\end{corollary}

\begin{proof}[Proof sketch]
Every member of every slice lies in $\Par$ by Definition~\ref{def:slice};
conversely, Lemma~\ref{lem:coverage} places each $(j,n)\in\Par$ in the slice
of its canonical circuit count.
\end{proof}

\section{The uniform-over-\texorpdfstring{$R$}{R} connection}\label{sec:uniform}

The formal development accompanying \cite{Roz25} proves, for each fixed
$R\ge 1$, that the slice $\ParR{R}$ is finite, conditionally on an excursion
estimate of polynomial type imposed along the members of the slice --- a
windowed form of the bound $M_T(n)\le n^{\beta}$ of \cite[eq.~(16)]{Roz25},
which is derived there from Conjecture~6.2.  Internally, that
per-$R$ theorem derives a bound $j\le J(R)$ on the lengths occurring in
$\ParR{R}$.  It is natural to ask what a single bound $J$ valid for
\emph{all} $R$ would give.  The answer is: everything.

\begin{theorem}\label{thm:uniform}
The following are equivalent.
\begin{enumerate}
\item[(a)] There is a $J\in\N$ such that $j\le J$ for every $R$ and every
$(j,n)\in\ParR{R}$.
\item[(b)] The paradoxical set $\Par$ is finite.
\end{enumerate}
\end{theorem}

\begin{proof}[Proof sketch]
(a)$\Rightarrow$(b): by Lemma~\ref{lem:coverage} every $(j,n)\in\Par$ lies in
some slice, so $j\le J$ holds on all of $\Par$, and
Proposition~\ref{prop:lengthbound} yields finiteness.
(b)$\Rightarrow$(a): if $\Par$ is finite, the set of lengths occurring in
$\Par$ is finite, hence bounded by some $J$; and every member of every slice
belongs to $\Par$, so $j\le J$ on every slice.
\end{proof}

\begin{remark}\label{rem:collatz}
By \cite[Cor.~3.3]{Roz25}, the finiteness of the paradoxical set implies the
Collatz conjecture.  Theorem~\ref{thm:uniform} therefore makes precise the
observation that a length bound uniform over all $R$ ``would be Collatz
itself'': the uniform-over-$R$ version of the per-$R$ finiteness theorem is
not merely sufficient for, but equivalent to, the finiteness of the
paradoxical set.
\end{remark}

\begin{remark}\label{rem:notclaimed}
We emphasize what is \emph{not} claimed: assuming the excursion estimate of
\cite[Conj.~6.2]{Roz25} uniformly in $R$ does not produce the uniform length
bound of Theorem~\ref{thm:uniform}(a) through the per-$R$ argument.  That
argument extracts a collision between circuits via a pigeonhole exponent of
order $j/(2R)$; for a sequence with $R\approx j/2$ circuits this exponent
degenerates to a constant and the collision disappears.  This gap is exactly
the Collatz-hardness of the uniform statement.
\end{remark}

\appendix

\section{Formal statements}\label{app:lean}

The results of this note are machine-verified in Lean~4, in the file
\texttt{paradoxical/CanonicalShape.lean} of the author's formal development;
all proofs are complete and free of extra-logical axioms.  The table below
maps each environment to its formal counterpart.  Statements recalled from
companion files are marked accordingly.

\begin{center}
\begin{tabular}{lll}
\toprule
Environment & Lean declaration & status \\
\midrule
Def.~\ref{def:paradoxical} & \texttt{RT.IsParadoxical} & definition (companion file) \\
Def.~\ref{def:slice} & \texttt{RCircuitSlice} & definition (companion file) \\
Prop.~\ref{prop:lengthbound} & \texttt{RT.finitely\_many\_paradoxical\_of\_} & proved (companion file) \\
 & \quad\texttt{paradoxical\_length\_bounded} & \\
Def.~\ref{def:shape} & \texttt{Paradoxical.shapeOf} & definition \\
Prop.~\ref{prop:roundtrip} & \texttt{Paradoxical.wordOfShape\_shapeOf} & proved \\
Def.~\ref{def:circuitcount} & \texttt{Paradoxical.circuitCount} & definition \\
Lem.~\ref{lem:coverage} & \texttt{Paradoxical.rCircuitSlice\_circuitCount} & proved \\
Cor.~\ref{cor:exhaustion} & \texttt{Paradoxical.iUnion\_rCircuitSlice} & proved \\
Thm.~\ref{thm:uniform} & \texttt{Paradoxical.uniform\_slice\_length\_} & proved \\
 & \quad\texttt{bound\_iff\_finitely\_many\_paradoxical} & \\
 & (forward direction also separately: & \\
 & \quad\texttt{Paradoxical.finitely\_many\_paradoxical\_} & \\
 & \quad\texttt{of\_uniform\_slice\_length\_bound}) & \\
\bottomrule
\end{tabular}
\end{center}

% omitted from the exposition (folded into prose or definitional):
%   Paradoxical.shapeOf_ne_nil      -- nonemptiness of the decomposition
%   Paradoxical.wlen_shapeOf        -- length additivity (stated inside Prop. 2.2)
%   Paradoxical.length_pbits        -- |V_j(n)| = j (definitional)
%   Paradoxical.pbits_binary        -- V_j(n) is a binary word (definitional)
%   Paradoxical.circuitCount_pos    -- R_j(n) >= 1 for j >= 1 (stated in prose)

\begin{thebibliography}{9}

\bibitem{Roz25}
O.~Rozier and C.~Terracol,
\emph{Paradoxical behavior in Collatz sequences},
arXiv:2502.00948 (2025).

\end{thebibliography}

\end{document}
